% Generated draft: requires mathematical and source review.
\documentclass[11pt]{article}
\usepackage[margin=25mm]{geometry}
\usepackage{fontspec}
\setmainfont{Latin Modern Roman}
\usepackage{amsmath,amssymb,mathtools,unicode-math}
\setmathfont{Latin Modern Math}
\usepackage{enumitem,needspace,xurl,hyperref}
\setlength{\parindent}{0pt}
\setlength{\parskip}{5pt}
\setlength{\emergencystretch}{3em}
\newcommand{\N}{\mathbb{N}}
\newcommand{\Z}{\mathbb{Z}}
\newcommand{\Q}{\mathbb{Q}}
\newcommand{\R}{\mathbb{R}}
\newcommand{\C}{\mathbb{C}}
\newcommand{\F}{\mathbb{F}}
\newcommand{\A}{\mathbb{A}}
\newcommand{\PP}{\mathbb{P}}
\newcommand{\E}{\mathbb{E}}
\newcommand{\eps}{\varepsilon}
\newcommand{\dd}{\,\mathrm{d}}
\begin{document}
\section*{UnsolvedMath ID: 20000412}
\textbf{Source problem code:} catalogue entry
\subsection{Definitions and mathematical statement}
For a variety $X$ over a number field $k$, a Tate class in a chosen degree and twist means a class in $H^{2r}_{\mathrm{et}}(X_{\overline k},\Q_\ell(r))$ fixed after some finite extension of $k$; equivalently one considers the union of invariant spaces for open Galois subgroups. For a nonproper Hilbert modular variety one must specify ordinary, interior, intersection, or compactification cohomology and boundary contributions. The source does not do so.

\textbf{Original Problem 3 (J. Getz).} Let $X$ be a Hilbert modular variety, whose cohomology is described heuristically as resembling a product of modular-curve cohomologies. Under what circumstances does it have Tate classes not coming from modular subvarieties? In the automorphic submotive or Galois-representation interpretation, must these exotic classes arise from CM automorphic forms, as in the workshop's Question 2? The relevant modular subvarieties, their cycle-class span, and the content of Question 2 are not supplied. Thus ``exotic'' requires a specified comparison subspace, not merely the existence of Tate classes.

\textbf{Attached unreviewed literature account.} The aid attributes to Getz--Hahn results for a cohomological nondihedral eigensystem on a Galois Hilbert modular variety, with sufficiently large $\ell$. It says the eventual Tate Hecke rank is zero when the projective-conjugacy stabilizer has odd order $m$, and a Catalan power when $m$ is even; a Tate line forces quadratic base-change descent after twisting; and rank-at-most-one Tate spaces are generated by twisted Hirzebruch--Zagier cycles. The summary gives no exact rank exponent, stabilizer definition, field and level hypotheses, or publication locator, so this is not a complete theorem statement. It does not identify every nonmodular class with a CM automorphic form.

The aid separately proposes a character-valued exact-base-change lemma. Let $L/F$ be Galois of degree $2n$, let $\rho_0$ be a two-dimensional characteristic-zero Galois representation in the intended full-$\mathrm{SL}_2$ monodromy setting, and form the tensor-induced representation $\operatorname{As}_{L/F}(\rho_0|_{G_L})$. After the full determinant normalization, represented in the tensor model by twisting by $(\det\rho_0)^{-n}$, the aid identifies its Tate multiplicity space with
\[\operatorname{Sp}^{(n,n)}\big|_{\operatorname{Gal}(L/F)\hookrightarrow S_{2n}},\]
where $\operatorname{Sp}^{(n,n)}$ is the characteristic-zero Specht module and the embedding is the regular permutation action on the $2n$ tensor factors. Its dimension is the Catalan number $\binom{2n}{n}/(n+1)$ in this tensor model. Precise monodromy, tensor-induction sign, Galois and normalization hypotheses require review.

For degree four, the stated restrictions are
\[\operatorname{Sp}^{(2,2)}|_{C_4}\cong\mathbf1\oplus\chi_2,\qquad\operatorname{Sp}^{(2,2)}|_{V_4}\cong\mathbf1^{\oplus2},\]
where $\chi_2$ is the unique nontrivial quadratic character of $C_4$. This refines rank to fixed-character multiplicities. The aid expressly constructs no new algebraic cycle and gives no general CM classification; representation-theoretic multiplicities alone do not supply those missing conclusions.
\subsection{Short English statement}
Getz asks when Hilbert modular varieties have Tate classes beyond modular-subvariety classes, and whether CM forms explain them. The aid reports qualified Getz--Hahn results and an exact-base-change Specht-module refinement, including the different cyclic and Klein-four character decompositions. It supplies neither new cycles nor a general answer to the CM question.
\subsection{Sources}
\begin{itemize}
\item Local UnsolvedMath registry, permanent ID 20000412, source code catalogue entry; \texttt{aim-algebraic-number-theory-notes.json}, notes[63], tag \texttt{problem}.
\item \emph{Problems from AIM Tate Conjecture Workshop}, July 23--27, 2007, Problem 3 (J. Getz): \url{https://aimath.org/WWN/tateconjecture/tateconjecture.pdf}. The reference to Question 2 is retained but its content is not supplied. The verification says the source was inspected; this draft claims no new PDF visit.
\item Attached machine-generated research attempt 1 names published Getz--Hahn results without bibliographic locators. Result \texttt{partial\_result}, status \texttt{partially\_solved}, rigor label \texttt{proved}, novelty confidence low. It is not peer reviewed and requires independent verification before citation. Original literature field empty.
\end{itemize}
\end{document}
